\exercise{Flight response\label{exStripesFlight}}{3}
Consider again the predator-prey system from the previous exercise. But now the predator, $y$, does not diffuse, while the prey $x$ tries to avoid the predators. Assume that the departure rate of prey individuals from a patch is $axy$, where $a$ is a parameter. 

\subquestion Consider the system on the linked-pair network. Write the differential equations and show that the steady state from the previous exercise is still stationary.

\solution 
We consider two patches, 1 and 2, connected by a link of weight one. The differential equations for the two patch system are 
\eqa{
\dot{x}_1&=& (4{x_1}^2+{x_1})/5 -x_1y_1 - ax_1y_1 + ax_2y_2 \\
\dot{y}_1&=& x_1y_1 - {y_1}^2 \\
\dot{x}_2&=& (4{x_2}^2+x_2)/5 -x_2y_2 - ax_2y_2 + ax_1y_1 \\
\dot{y}_2&=& x_2y_2 - {y_2}^2 
}
If we substitute in our previous steady state $x=y=1$ we find 
\eqa{
\dot{x}_1&=& (4+1)/5 -1 - a + a = 0 \\
\dot{y}_1&=& 1 - 1 = 0\\
\dot{x}_2&=& (4+1)/5 - 1 - a + a = 0\\
\dot{y}_2&=& 1 - 1 = 0
}
so our previous steady state is still a steady state. 

\subquestion 
Show that the Jacobian for the linked pair can still be written in the form ${\bf J}={\bf I}\otimes{\bf P}-{\bf L}\otimes {\bf C}$, albeit with a different matrix $\bf C$. 

\solution
The Jacobian is 
\eqa{
{\bf J} = \avecccc{
4/5-a & -1 -a & a      & a \\
1     & -1    & 0      & 0 \\
a     & a     & 4/5-a  & -1-a \\ 
0     & 0     & 1      & -1
}
}
Compare this with the Jacobian of a single isolated note (where there is nowhere to go and hence also no departures) 
\eq{
{\bf P} = \avecc{4/5 & -1 \\ 1 & -1 }.
}
We can see now that the single-node Jacobian still appears in blocks on the diagonal. Hence we can write tht Jacobian of the linked pair as 
\eq{
{\bf J} = \avecc{{\bf P} & 0 \\ 0 & {\bf P}} + \avecccc{
-a & -a &  a & a \\
0  &  0 &  0 & 0 \\
a  &  a & -a & -a \\ 
0  &  0 &  0 & 0
}
}
We can now see that the rest in the second term prominently contains the block 
\eq{
{\bf C} = \avecc{a & a \\ 0 & 0}
}
this is the new coupling matrix. Using it the Jacobian can be written as 
\eq{
{\bf J} = \avecc{{\bf P} & 0 \\ 0 & {\bf P}} + \avecc{ -{\bf C} & {\bf C} \\ {\bf C} & -{\bf C}}
}
or equivalently 
\eq{
{\bf J} = {\bf I} \otimes {\bf P} - {\bf L} \otimes {\bf C}
}
where 
\eq{
{\bf L} = \avecc{1 & -1 \\ -1 & 1 }
}
is the Laplacian of the unweighted linked pair. 

\subquestion
Is the homogeneous state stable in this case?

\solution 
Because we have the structure ${\bf J}={\bf I}\otimes{\bf P}-{\bf L}\otimes {\bf C}$ we know that we can compute the eigenvalues from 
${\bf P}-\kappa {\bf C}$. We could now substitute the eigenvalues
$\kappa_1=0$ and $\kappa_2=2$ for the linked pair, but using our master stability function trick we can as well explore the general case for arbitrary $\kappa$. 

The relevant matrix in this case is 
\eq{
{\bf P}-\kappa {\bf C} = \avecc{4/5-\kappa a & -1 - \kappa a \\ 1 & -1} 
}
Note that $\kappa$ and $a$ appear only together as a product, so only their product can matter and specifying either of them would not save us any work.

We already know from the previous exercise that at $\kappa a =0$ the system is stable. So instead of solving another quadratic polynomial to compute the eigenvalues, let's just check if there are any bifurcations in which the stability is lost. 

Hopf like bifurcations can only occur if the trace of the Jacobian vanishes, but in this case
\eq{
{\rm Tr} ({\bf P}-\kappa {\bf C}) = 4/5 -1 -\kappa a = -1/5 -\kappa a
}
We know $\kappa\geq 0$ (Laplacians have non-negative eigenvalues) and $a>0$ (a negative departure rate would not make sense). Therefore $a\kappa\geq 0$, the trace is always negative and Hopf-like bifurcations cannot occur. 

To check for fold-like bifurcations we check the determinant
\eq{
0 = -(4/5-\kappa a) + (1+\kappa a) = 1/5 
}
so there are no fold-like bifurcations either. Hence this sort of flight response won't destabilize the stationary state. However, generally cross-diffusion (i.e.~diffusion where the departure rate of one variable depends on other variables) is likely to cause diffusion-driven instabilities.